% !TEX program = xelatex
% !TeX encoding = UTF-8
\documentclass[12pt,a4paper]{article}
\usepackage{fontspec}
\usepackage{polyglossia}
\setmainlanguage{japanese}
\setotherlanguage{english}

% ============================================================
% 日本語フォント – システムにインストールされている Meiryo を使用 (Windows)
% Script=CJK が正しい指定 (Japanese は認識されない)
% ============================================================
\setmainfont{Meiryo}[Script=CJK, Language=Japanese]
\setsansfont{Meiryo}[Script=CJK, Language=Japanese]
\setmonofont{MS Gothic}[Script=CJK, Language=Japanese]

% polyglossia 用の等幅日本語フォント（同じく MS Gothic）
\newfontfamily\japanesefonttt{MS Gothic}[Script=CJK, Language=Japanese]

% 英数字フォント (数式や参考文献用)
\newfontfamily\englishfont{Times New Roman}[Ligatures=TeX]

% ============================================================
% 数学・レイアウトパッケージ
% ============================================================
\usepackage{amsmath,amssymb,amsthm,mathtools}
\usepackage{geometry}
\geometry{left=2.5cm,right=2.5cm,top=2.5cm,bottom=2.5cm}
\usepackage{booktabs}
\usepackage{hyperref}
\usepackage{url}
\usepackage{enumitem}
\usepackage{float}
\usepackage{xcolor}
\usepackage{titlesec}
\usepackage{fancyhdr}

% 色定義
\definecolor{gold}{RGB}{255,215,0}
\definecolor{darkblue}{RGB}{0,0,139}
\definecolor{skyblue}{RGB}{0,102,204}

% YUCT マクロ
\newcommand{\Keff}{K_{\mathrm{eff}}}
\newcommand{\kappac}{\kappa_c}
\newcommand{\eps}{\varepsilon}
\newcommand{\df}{d_{\!f}}
\newcommand{\dint}{\ensuremath{D\!+\!I\!\cdot\!R}}
\newcommand{\Sodd}{S_{\mathrm{odd}}}
\newcommand{\Seven}{S_{\mathrm{even}}}

% 日本語定理環境
\newtheorem{theorem}{定理}[section]
\newtheorem{definition}[theorem]{定義}
\newtheorem{proposition}[theorem]{命題}
\newtheorem{conjecture}[theorem]{予想}
\theoremstyle{remark}
\newtheorem{remark}[theorem]{備考}

\hypersetup{
	colorlinks=true,
	linkcolor=blue,
	citecolor=blue,
	urlcolor=blue,
}

% 見出し書式
\titleformat{\section}{\LARGE\bfseries\color{darkblue}}{\arabic{section}}{1em}{}
\titleformat{\subsection}{\Large\bfseries\color{darkblue}}{\arabic{subsection}}{1em}{}

% フッター
\fancypagestyle{firstpage}{%
	\fancyhf{}%
	\renewcommand{\headrulewidth}{0pt}%
	\renewcommand{\footrulewidth}{0pt}%
	\fancyfoot[C]{%
		\begin{minipage}[t]{\textwidth}
			\centering
			\textcolor{gold}{\rule{\textwidth}{1.6pt}}\\[5pt]
			{\small \textsuperscript{1}Yakushev Research, \url{https://yuct.org/}} \hfill
			{\small YPSDC プロトコル: \url{https://ypsdc.com/}}\\[-2pt]
			\textcolor{gold}{\rule{\textwidth}{1.6pt}}\\[5pt]
			\textbf{Yakushev 統一調整理論 2026} \hfill \thepage\\[10pt]
		\end{minipage}%
	}%
}

\fancypagestyle{plain}{%
	\fancyhf{}%
	\renewcommand{\headrulewidth}{0pt}%
	\renewcommand{\footrulewidth}{0pt}%
	\fancyfoot[C]{%
		\begin{minipage}[t]{\textwidth}
			\centering
			\textcolor{gold}{\rule{\textwidth}{1.6pt}}\\[4pt]
			\textbf{Yakushev 統一調整理論 2026} \hfill \thepage\\[4pt]
		\end{minipage}%
	}%
}

\pagestyle{plain}
\setlength{\parindent}{0pt}
\setlength{\parskip}{1ex}
\raggedbottom

\makeatletter
\let\@ensure@LTR\relax
\makeatother

% ============================================================
% 文書開始
% ============================================================
\begin{document}
	
	\begin{titlepage}
		\centering
		\vspace*{1cm}
		
		\Huge\textbf{付録 F（改訂版）: 経済辞書のフラクタルダイナミクスと貨幣の調整理論}
		
		\vspace{0.5cm}
		\Large\textbf{為替レート、通貨圏、自己組織化的経済境界の理論的枠組み}
		
		\vspace{1.5cm}
		
		\Large\textbf{アレクセイ・V・ヤクシェフ\textsuperscript{1}}
		
		YUCT 理論: \url{https://yuct.org/}\\
		YPSDC プロトコル: \url{https://ypsdc.com/}
		
		\vspace{2cm}
		\textbf{DOI:} \url{https://doi.org/10.5281/zenodo.18444598}
		
		\vspace{1cm}
		
		\large
		2026年6月 \\ \large バージョン 3.0（旧付録 F を置き換える）
		
		\vspace{2cm}
		
		\begin{minipage}{0.9\textwidth}
			\begin{abstract}
				\noindent
				本付録は、YUCT の経済学分野へのアプローチを根本的に見直すものである。我々は、GDP や市場ボラティリティを普遍指数 \(\beta = 0.67\) に当てはめようとする過去のすべての試みを放棄し、代わりに調整オントロジーに基づく第一原理からの理論的枠組みを構築する。経済は相互作用する \textbf{経済辞書}（企業、セクター、国家）のネットワークとしてモデル化され、各辞書は調整効率 \(K_{\mathrm{eff}}\) によって特徴付けられる。任意の二つの辞書間の境界は、マンデルブロ集合に類似したフラクタルダイナミクスを示すことを示す：自己相似性、複雑性、そして普遍誤差法則 \(\varepsilon = \kappa_c \alpha K_{\mathrm{eff}}^{-\beta}\) によって支配される。貨幣は辞書境界を越えた活性化を調整する \textbf{調整パラメータ} として定義され、従って為替レートは調整効率の比によって決定される。この枠組みは、いくつかの定量的で反証可能な予測を提供する：(1) 高頻度為替レート時系列のハースト指数は長い時間スケールで \(H \approx 1/\beta \approx 1.5\) に収束する；(2) 金融収益率のマルチフラクタルスペクトルは、一次代数サイクルによって決定される普遍的な左歪性を示す；(3) 為替レートのボラティリティ \(\sigma\) は、二つの辞書間の調整ギャップ \(\Delta K_{\mathrm{eff}}\) に対して \(\sigma \propto (\Delta K_{\mathrm{eff}})^{-0.67}\) でスケーリングする。本付録では実証的フィッティングは一切行わず、すべての予測は経済システムにおいて \(K_{\mathrm{eff}}\) の信頼できるプロキシが得られ次第テスト可能な形で定式化される。この作業は、YUCT 経済学を疑わしい相関の寄せ集めから、厳密で反証可能な研究プログラムへと転換する。
			\end{abstract}
		\end{minipage}
		
		\vspace{1cm}
		
		\noindent
		\small\textbf{キーワード:} YUCT, 調整効率, 経済辞書, 為替レート, フラクタル境界, マンデルブロ集合, ハースト指数, マルチフラクタル解析, 調整パラメータとしての貨幣.
		
		\vspace{0.5cm}
		
		\noindent
		\textcopyright\ 2026 Yakushev Research. All rights reserved.
	\end{titlepage}
	
	\tableofcontents
	\newpage
	
	% ============================================================
	% 第1章: はじめに – なぜ新たな出発が必要か
	% ============================================================
	\section{はじめに: なぜ新たな出発が必要か}
	
	本付録の以前のバージョンでは、マクロ経済時系列（GDP 成長ボラティリティ、暗号通貨予測誤差、地政学的リスク指標など）が YUCT の普遍スケーリング則 \(\varepsilon \propto K_{\mathrm{eff}}^{-0.67}\) に従うことを実証しようとした。その後の独立した検証により、言及された多くの表や回帰分析が誤っているか、事後的に好都合なケースを選択したものであることが明らかになった。本改訂版はこれらの欠陥を完全に認め、過去のすべての実証的主張を撤回する。
	
	根本的な誤りは YUCT フレームワーク自体にあるのではなく、「経済辞書」が実際に何であり、その調整効率をどのように測定できるかについての適切な理論モデルを事前に開発せずに、直接的な実証的フィッティングを強要しようとした点にある。経済学は素粒子物理学ではない。経済主体は点粒子ではなく、その質量は基本定数ではない。したがって、経済データから単純な質量スケールを探すことは当初から誤った方向であった。
	
	本付録は根本的に異なるアプローチをとる。経済を、そのサイズが離散スペクトルに属さねばならない対象の集合として扱う代わりに、我々はそれを相互作用する \textbf{辞書の動的ネットワーク} と見なし、その境界が興味深い現象の主要な場であるとする。為替レート、金利、金融ボラティリティは孤立した主体の特性ではなく、辞書間の \textbf{境界の創発的特性} である。我々は、これらの境界の幾何学がフラクタルかつ自己相似的であり、YUCT の他のすべての領域に現れる同じ普遍指数 \(\beta = 0.67\) によって支配されると主張する。
	
	本付録では実証的フィッティングは一切行わない。すべての定量的主張は YUCT の公理から数学的に導出され、将来の研究のための反証可能な予測として提示される。
	
	% ============================================================
	% 第2章: 経済辞書とその境界
	% ============================================================
	\section{経済辞書とその境界}
	
	\subsection{経済辞書の定義}
	
	\begin{definition}[経済辞書]
		\textbf{経済辞書} \(\mathcal{D}_i\) は、経済主体（企業、家計、セクター、国家）が外部指標に応答して活性化できる可能な行動、契約、生産プロトコル、価格ルール、制度規範の構造化された集合である。辞書のサイズはシャノンエントロピー \(H(\mathcal{D}_i)\) で測定される。
	\end{definition}
	
	例:
	\begin{itemize}
		\item 企業の辞書には、製品カタログ、供給契約、価格設定アルゴリズム、内部管理手順が含まれる。
		\item 国家経済の辞書には、法体系、税法、貿易協定、金融政策の枠組み、金融インフラが含まれる。
	\end{itemize}
	
	辞書 \(i\) の調整効率は
	\begin{equation}
		K_{\mathrm{eff}}^{(i)} = \frac{H(\mathcal{A}_i)}{H(\kappa_i)},
		\label{eq:keff_econ}
	\end{equation}
	ここで \(\mathcal{A}_i\) は辞書が生成できる行動の集合、\(\kappa_i\) はそれらを引き起こす指標（価格シグナル、規制指示、市場注文）である。
	
	\subsection{二つの辞書間の境界}
	
	\begin{definition}[辞書境界]
		二つの経済辞書 \(\mathcal{D}_i\) と \(\mathcal{D}_j\) の間の \textbf{境界} \(\partial\mathcal{D}_{ij}\) は、一方の辞書で発生した指標が他方の辞書の入力を活性化するすべての取引、契約、情報交換の集合である。
	\end{definition}
	
	境界は以下の現象の自然な場である:
	\begin{itemize}
		\item 為替レート（国家辞書の場合）、
		\item 取引コストと価格差（企業レベルの辞書の場合）、
		\item 情報の非対称性と逆選択。
	\end{itemize}
	
	\begin{proposition}[境界における調整パラメータ]
		境界 \(\partial\mathcal{D}_{ij}\) を介した調整効率は、無次元の \textbf{調整パラメータ} \(\chi_{ij}\) によって測定され、次のように定義される:
		\begin{equation}
			\chi_{ij} = \frac{K_{\mathrm{eff}}^{(j)}}{K_{\mathrm{eff}}^{(i)}} \cdot \rho_{ij},
			\label{eq:chi}
		\end{equation}
		ここで \(\rho_{ij}\) は二つの辞書の互換性（共通の標準、共通言語、法的調和など）を決定する共鳴因子である。
	\end{proposition}
	
	\textbf{調整パラメータとしての貨幣。} 独立した通貨を持つ二つの国家経済の場合、調整パラメータ \(\chi_{ij}\) は正規化定数を除いて正確に \textbf{為替レート}（通貨 \(i\) 当たりの通貨 \(j\) の単位）である。\(K_{\mathrm{eff}}^{(i)}\) に対する \(K_{\mathrm{eff}}^{(j)}\) の値が大きいことは、辞書 \(j\) がより「価値がある」ことを意味し — その入力がより確実に活性化される — したがってその通貨はプレミアムを要求する。
	
	% ============================================================
	% 第3章: 境界のフラクタルダイナミクス
	% ============================================================
	\section{境界のフラクタルダイナミクス}
	
	\subsection{マンデルブロ集合との類似}
	
	マンデルブロ集合 \(\mathcal{M}\) は反復写像 \(z_{n+1} = z_n^2 + c\)（\(z_0 = 0\)）で定義される。数列が有界に留まる点 \(c\) は \(\mathcal{M}\) に属し、発散する点はその外側にある。境界 \(\partial\mathcal{M}\) はフラクタルであり、無限に複雑で、すべてのスケールで自己相似性を示す。
	
	我々は経済的類似を提案する。二つの辞書 \(\mathcal{D}_i\) と \(\mathcal{D}_j\) が離散時間ステップ \(n\) でインデックス付けられた一連の取引を通じて相互作用するとする。\(n\) 回の取引後の調整状態は複素変数 \(z_n \in \mathbb{C}\) で記述され、その実部と虚部は二つの直交次元（例えば価格同期と契約履行）における類似度を表す。更新規則は
	\begin{equation}
		z_{n+1} = z_n^2 + \chi_{ij},
		\label{eq:mandel_econ}
	\end{equation}
	ここで \(\chi_{ij} \in \mathbb{C}\) は式~(\ref{eq:chi}) の（複素）調整パラメータである。絶対値 \(|\chi_{ij}|\) は結合強度を測定し、位相は二つの辞書の活性化周期間の位相ずれを測定する。
	
	\begin{conjecture}[マンデルブロの経済的対応]
		辞書のペア \((\mathcal{D}_i, \mathcal{D}_j)\) は、\(\chi_{ij}\) が一般化マンデルブロ集合 \(\mathcal{M}_{\mathrm{econ}}\) 内にある場合にのみ有界で安定な調整を維持できる（すなわち数列 \(\{z_n\}\) が有界に留まる）。境界 \(\partial\mathcal{M}_{\mathrm{econ}}\) は安定な調整領域（安定な為替レート、予測可能な貿易フロー）とカオス的発散領域（通貨危機、超インフレーション、貿易崩壊）を分離する。この境界の幾何学はフラクタルであり、すべての時間スケールで自己相似構造を持つ。
	\end{conjecture}
	
	\subsection{境界における普遍スケーリング}
	
	境界上のダイナミクスは YUCT 全体に適用される同一の普遍調整誤差法則に従うため、境界近傍の為替レートの統計的性質は普遍スケーリング則に従うはずである。\(\sigma_{ij}(\Delta t)\) を時間軸 \(\Delta t\) で測定された辞書 \(i\) と \(j\) の間の為替レートのボラティリティとする。すると普遍誤差法則 \(\varepsilon \propto K_{\mathrm{eff}}^{-\beta}\) から次を得る:
	
	\begin{proposition}[為替レートボラティリティのスケーリング]
		調整パラメータ \(\chi_{ij}\) のボラティリティは、調整ギャップ \(\Delta K_{\mathrm{eff}} = |K_{\mathrm{eff}}^{(i)} - K_{\mathrm{eff}}^{(j)}|\) に対して次のようにスケーリングする:
		\begin{equation}
			\sigma(\chi_{ij}) \propto (\Delta K_{\mathrm{eff}})^{-\beta} = (\Delta K_{\mathrm{eff}})^{-0.67}.
			\label{eq:vol_scaling}
		\end{equation}
		非常に類似した調整効率を持つ辞書のペア（小さなギャップ \(\Delta K_{\mathrm{eff}}\)）は高いカオス的ボラティリティを示し、大きなギャップを持つペアは低く安定したボラティリティを示す。
	\end{proposition}
	
	これは反証可能な予測である。\(K_{\mathrm{eff}}^{(i)}\) の信頼できるプロキシ（例えば世界銀行のガバナンス指標と貿易複雑性データの組み合わせ）が開発されれば、式~(\ref{eq:vol_scaling}) は国別パネルデータでテストできる。
	
	\subsection{金融時系列のフラクタル特性}
	エコノフィジックスでよく知られた経験的事実として、金融時系列はマルチフラクタル特性（ボラティリティクラスタリング、長期記憶、太い裾の収益率分布）を示す。YUCT 内では、これらは偶然ではなく、辞書境界のフラクタルダイナミクスの必然的な結果である。
	
	\begin{proposition}[境界におけるハースト指数]
		安定調整の臨界境界近傍で動作する辞書ペアについて、為替レート時系列のハースト指数 \(H\) は
		\begin{equation}
			H \to \frac{1}{\beta} = \frac{3}{2} \quad \text{as} \quad \Delta t \to \infty.
			\label{eq:hurst}
		\end{equation}
		より短い時間スケールでは、\(H\) はより低い値へのクロスオーバーを示し、これは境界のマルチフラクタル性を反映する。
	\end{proposition}
	
	この予測は、高頻度為替レートデータに対する標準的な R/S 解析または DFA によってテストできる。\(H = 1.5\) という値は、非常に長い記憶を持つプロセスの特徴であり、特異点に近い — これは大規模金融危機（例えば 2008 年危機）の直前にまさに観察される挙動である。
	
	\subsection{単一辞書の内部フラクタル組織}
	
	辞書間の境界がフラクタルかつカオス的であるのに対し、よく調整された辞書の内部構造は離散的な階層的秩序を示す。辞書の項目は一様に分布せず、規則、契約、プロトコルの自己相似分岐ツリーを形成する。企業規模、所得分布、都市システムの階層構造はすべてこの内部フラクタル構造の現れである。
	
	以前の失敗した試み（経済全体を質量スケールに当てはめようとした）とは異なり、我々は今や \textbf{一つの} 辞書内部において、下位単位（企業、都市、所得階層）のサイズがフラクタル次元 \(d_f\) および普遍的な \(\beta\) に関連した指数を持つべき乗則に従うと予測する。都市に関するよく知られたジップの法則（指数 \(\approx 1\)）や所得に関するパレートの法則（指数 \(\approx 1.5\)）は孤立した奇異ではなく、同一の基礎的調整幾何学の異なる投影である。
	
	% ============================================================
	% 第4章: 数学的枠組み
	% ============================================================
	\section{数学的枠組み: 境界における調整のラグランジアン}
	
	上述の定性的描像を数学的に精密化するため、我々はそれを 19 次元 YUCT ラグランジュ形式に埋め込む。経済セクターを場 \(\Psi_{72}(x^\mu, y^m)\) で記述する。ここで \(x^\mu\) は 4 つの時空座標、\(y^m\) は 15 個の内部（コンパクト）座標である。辞書 \(\mathcal{D}_i\) は特定の真空期待値 \(\langle \Psi_{72} \rangle_i\) に対応し、二つの辞書間の境界は \(\langle \Psi_{72} \rangle_i\) と \(\langle \Psi_{72} \rangle_j\) の間を内挿する領域壁解である。
	
	境界上の有効作用は
	\begin{equation}
		S_{\mathrm{boundary}} = \int d^4x \sqrt{-g} \left[ \frac{1}{2} \kappa (\nabla \chi)^2 + V(\chi) \right],
		\label{eq:action}
	\end{equation}
	ここで \(\chi(x)\) は局所調整パラメータ（為替レート場）、\(\kappa\) は境界剛性、ポテンシャル \(V(\chi)\) は
	\begin{equation}
		V(\chi) = V_0 \left( e^{\lambda(\chi - \chi_0)} - \lambda(\chi - \chi_0) - 1 \right), \qquad \lambda = \frac{1}{\beta} = \frac{3}{2}.
		\label{eq:potential}
	\end{equation}
	これは重力セクター（付録 G）でスカラー調整場を支配するものと同じポテンシャルであり、普遍性を保証する。
	
	\(\chi\) の古典的運動方程式は
	\begin{equation}
		\kappa \Box \chi - V'(\chi) = 0,
		\label{eq:eom}
	\end{equation}
	であり、その統計的変動は
	\begin{equation}
		\langle (\delta \chi)^2 \rangle \propto \int \frac{d^4k}{(2\pi)^4} \frac{1}{\kappa k^2 + V''(\chi_0)} \propto \chi_0^{-2\beta},
		\label{eq:fluctuations}
	\end{equation}
	を満たし、式~(\ref{eq:vol_scaling}) のボラティリティスケーリング則を再現する。
	
	% ============================================================
	% 第5章: 反証可能な予測
	% ============================================================
	\section{反証可能な予測}
	
	上記で開発された枠組みは、データフィッティングなしでテスト可能ないくつかの定量的予測を提供する:
	
	\begin{enumerate}[label=(\arabic*)]
		\item \textbf{為替レートのハースト指数。} 主要な変動相場制通貨ペア（例: EUR/USD, USD/JPY）について、20 年間の日次データから計算された長期ハースト指数は \(1.3 \le H \le 1.7\) の範囲にあるべきである。主流の金融の帰無仮説（ランダムウォーク、\(H = 0.5\)）は高い信頼度で棄却されるべきである。
		
		\item \textbf{ボラティリティ–調整ギャップスケーリング。} \(K_{\mathrm{eff}}\) の信頼できるプロキシ（例えば世界銀行のガバナンス指標と貿易複雑性データの組み合わせ）が構築されれば、為替レートのクロスセクションボラティリティは絶対差 \(\Delta K_{\mathrm{eff}}\) に対して \(\sigma \propto (\Delta K_{\mathrm{eff}})^{-0.67}\) でスケーリングすべきである。指数は対数座標での線形回帰によって推定できる；予測は傾きが \(-0.67 \pm 0.10\) であることである。
		
		\item \textbf{マルチフラクタルスペクトルの非対称性。} 為替レート収益率の特異点スペクトル \(f(\alpha)\) は左歪性を示すべきであり、非対称係数 \(A_\alpha \approx 0.30 \pm 0.05\) を持つ。これは大きな調整下落が大きな調整上昇よりも優勢であることを反映する。これは公開データに対する標準的な MF-DFA プログラムでテストできる。
		
		\item \textbf{通貨危機前の臨界減速。} 通貨ペアが臨界境界に近づくとき（すなわち \(\Delta K_{\mathrm{eff}} \to 0\)）、ボラティリティの自己相関時間は発散すべきである。これは、主要通貨危機の前にハースト指数の着実な増加とマルチフラクタルスペクトルの収縮が観察されるべきであり、早期警告シグナルを提供することを予測する。
	\end{enumerate}
	
	\textbf{予測のステータス:} すべての予測は経済データへのフィッティングなしに YUCT 公理から数学的に導出される。これらは独立した研究者によって直接テスト可能な形で提示される。これらのうちの一つでも確認されれば、貨幣の調整理論に対する強力な証拠を提供するだろう；四つすべての反証は枠組みの抜本的修正を必要とするだろう。
	
	% ============================================================
	% 第6章: 未解決問題と今後の展望
	% ============================================================
	\section{未解決問題と今後の展望}
	
	\begin{enumerate}[label=(\roman*)]
		\item \textbf{国家 \(K_{\mathrm{eff}}\) の信頼できるプロキシの構築。} 最も緊急の課題は、公開データから国家経済の調整効率を推定する方法を開発することである。候補：法のエントロピー、輸出バスケットの複雑性（イダルゴ–ハウスマン指数）、制度の質（世界銀行ガバナンス指標）、市場ミクロ構造の情報測度。少数の参照国グループに対して較正された複合指数は、全パネルに外挿できる。
		
		\item \textbf{経済的マンデルブロ写像の数値シミュレーション。} 写像 \(z_{n+1} = z_n^2 + \chi\) を体系的に研究すべきである。どの \(\chi\) の値に対して数列は有界に留まるか？境界のフラクタル次元は何か？ノイズ（有限温度）を加えると構造にどのような影響を与えるか？これらの問題は簡単な Python スクリプトで扱うことができ、実際の為替レートダイナミクスとの予期せぬ定量的な関連性を明らかにする可能性がある。
		
		\item \textbf{確立されたエコノフィジックスとの連携。} 第5節に挙げた予測は、エコノフィジックスの既存の結果（金融収益率のマルチフラクタル性、ボラティリティの長期記憶、べき乗則の裾）とかなり重なる。この文献を YUCT のレンズで解釈した系統的レビューは、新たなデータ収集なしに即座に枠組みの検証を提供するだろう。
	\end{enumerate}
	
	% ============================================================
	% 第7章: 結論
	% ============================================================
	\section{結論}
	
	我々は YUCT の経済学へのアプローチを完全に再構築した。新しい枠組みは、経済的総量を質量スケールに強制するという失敗した戦略を放棄し、代わりに経済辞書とその境界に関する第一原理からの理論を発展させる。貨幣は辞書間の活性化を調整する調整パラメータとして定義され、為替レートは YUCT の他のすべての場所に現れる同じ普遍指数 \(\beta = 0.67\) によって支配されることが示される。結果として得られる予測は定量的で反証可能であり、事後のデータフィッティングとは無関係である。これにより、YUCT 経済学は負担から有望で厳格な研究プログラムへと変貌する。
	
	\section*{免責事項}
	
	本付録は理論的研究である。いかなる実証的主張も行わず、投資や政策の推奨を意図するものではない。すべての定量的予測は科学的審査のために提示され、将来のデータによって反証される可能性がある。
	
	\section*{データの入手可能性}
	
	この理論的研究のために新しいデータは生成されなかった。言及されたすべての情報源は公開されている。
	
	\begin{thebibliography}{99}
		\bibitem{AppendixA} A.V. Yakushev, ``付録 A: 19次元 YUCT ラグランジュ形式,'' in \emph{YUCT}, Zenodo (2026).
		\bibitem{AppendixL} A.V. Yakushev, ``付録 L: フラクタル調整誤差測定,'' in \emph{YUCT}, Zenodo (2026).
		\bibitem{AppendixG} A.V. Yakushev, ``付録 G: 幾何学的張力としての重力,'' in \emph{YUCT}, Zenodo (2026).
		\bibitem{AppendixX} A.V. Yakushev, ``付録 X: シャノン情報理論の一般化,'' in \emph{YUCT}, Zenodo (2026).
		\bibitem{Mandelbrot1982} B.B. Mandelbrot, \emph{自然界のフラクタル幾何学}, W.H. Freeman (1982).
		\bibitem{Pareto1896} V. Pareto, \emph{政治経済学講義}, Rouge (1896).
		\bibitem{Zipf1949} G.K. Zipf, \emph{人間行動と最小努力の原理}, Addison-Wesley (1949).
		\bibitem{Hidalgo2009} C.A. Hidalgo, R. Hausmann, ``経済的複雑性の構成要素,'' \emph{PNAS} \textbf{106}, 10570–10575 (2009).
		\bibitem{Kantelhardt2002} J.W. Kantelhardt 他, ``非定常時系列のマルチフラクタル・ディトレンド変動解析,'' \emph{Physica A} \textbf{316}, 87–114 (2002).
	\end{thebibliography}
	
\end{document}